\chapter*{怎样使用这本书}
\addcontentsline{toc}{chapter}{怎样使用这本书}
\markboth{怎样使用这本书}{怎样使用这本书}

第一次接触操作系统时，建议按目录顺序阅读。前四章解释后面会反复出现的
电路、逻辑和 CPU 概念；从第五章开始，每章都会回到项目中的文件、构建产物
和运行结果。

已经学过计算机组成的读者可以从第六章开始。遇到电源、时钟、逻辑电平或状态
存储等概念时，再回看前四章。无论从哪里进入，都建议亲手运行实验，不要只读
最终结论。

\section*{先观察，再打开代码}

学习一个阶段时，可以按下面的顺序进行：

\begin{enumerate}[leftmargin=2.2em]
  \item 运行该阶段的测试或 QEMU，用原样文本记录串口输出和退出状态。
  \item 找到产生这条输出的汇编或 C++ 函数，沿调用路径向前追踪。
  \item 手算其中一个地址、位字段、容量或状态变化。
  \item 改变一个输入，再比较输出走到了哪个分支。
\end{enumerate}

例如，UART 阶段先确认终端是否出现字符，再检查 Line Status Register 的
发送保持寄存器空位。长模式阶段先观察跳转前后的日志，再核对 CR0、CR4、
EFER 和页表。文件系统阶段先创建并读回一个文件，随后才分析 inode、位图和
磁盘写入顺序。

\section*{二十章会做什么}

\begin{osgridlongtable}{L{0.06\linewidth}L{0.28\linewidth}L{0.56\linewidth}}
章 & 主题 & 读完以后可以亲手完成的事情 \\ \hline
\endfirsthead
章 & 主题 & 读完以后可以亲手完成的事情 \\ \hline
\endhead
1 & 从测量和计数开始 &
带着单位计算时间、频率、容量和地址范围。\\ \hline
2 & 电路怎样传递能量和信号 &
阅读基本原理图，并计算电阻、电流、功率和 RC 时间。\\ \hline
3 & 机器怎样记住一个状态 &
从真值表画出组合逻辑，沿时钟边沿分析寄存器和状态机。\\ \hline
4 & CPU 怎样读到第一条指令 &
沿电源、复位、地址总线和 ROM 说明一次取指。\\ \hline
5 & 操作系统替程序做了什么 &
区分固件、Loader、Kernel、驱动和用户程序接手工作的时刻。\\ \hline
6 & x86 为什么从实模式醒来 &
手算复位地址、分段地址和 A20 回绕。\\ \hline
7 & CPU 怎样找到内存和设备 &
区分内存事务与 Port I/O，并在载板原理图中找到对应器件。\\ \hline
8 & 读懂 CPU 执行的指令 &
阅读 Intel 语法汇编，跟踪寄存器、RFLAGS 和内存变化。\\ \hline
9 & 把源码变成一块启动 ROM &
从源文件追到 ELF、链接布局、ROM 偏移和最终机器字节。\\ \hline
10 & 让机器第一次说话 &
初始化 UART，发送字符，并识别设备忙与超时。\\ \hline
11 & 从磁盘进入 64 位内核 &
读取 ATA 扇区、解析 ELF64，并检查进入长模式需要的寄存器。\\ \hline
12 & 内核怎样管理内存和异常 &
解读异常现场，手算页表索引，并跟踪页帧申请与释放。\\ \hline
13 & 中断怎样把设备带进内核 &
从设备 IRQ 追到 IDT 入口，再返回 Ring 3。\\ \hline
14 & 让多个线程轮流使用 CPU &
观察 Thread 状态变化、上下文切换和一次系统调用。\\ \hline
15 & 文件描述符怎样找到文件 &
从用户 fd 追到 FileDescription、Vnode、路径和挂载点。\\ \hline
16 & 进程怎样装入程序和分配内存 &
跟踪 exec、缺页、文件映射、fork 和写时复制。\\ \hline
17 & 文件怎样在断电后留下来 &
检查磁盘布局，并比较不同断电位置留下的文件系统状态。\\ \hline
18 & 进程怎样等待、通信和控制终端 &
运行 pipe、futex、signal 和前后台作业实验。\\ \hline
19 & 用户程序怎样使用系统服务 &
读取 ABI、devfs 和 procfs，并运行磁盘中的命令行工具。\\ \hline
20 & 回看整个系统并准备下一次修改 &
选择一个模块，找到接口、相关测试和修改后需要重跑的路径。\\ \hline
\end{osgridlongtable}

\section*{按问题选择阅读路线}

\topic{我想先看到系统启动。}
读第 4、6、8--13 章。第 4 章解释取指所需的硬件条件，第 6 章解释 x86
复位状态，第 8--11 章把源码变成 ROM 并进入内核，第 12--13 章加入保护和
用户态。

\topic{我想理解进程和虚拟内存。}
先读第 12--14 章，弄清异常、页表、中断和 Thread 切换，再读第 16 章。
fork 和写时复制都依赖这些基础。

\topic{我想理解文件和磁盘。}
读第 7、15、17 章。先区分设备端口和磁盘扇区，再沿 fd、VFS、inode、页
缓存和设备写入追踪一个字节。

\topic{我想理解 Shell 怎样控制程序。}
读第 13--19 章，重点观察系统调用、pipe、signal、TTY、进程组和前后台
切换。

\section*{遇到陌生概念时停一下}

看到新的缩写时，先确认它属于 CPU、设备、内核对象、用户接口还是磁盘格式。
看到地址时，先写出地址空间和单位。看到一个失败返回时，继续找调用者怎样
处理它。看到测试通过时，记录输入、输出和它没有覆盖的分支。

这本书不要求背下所有端口号和位值。更重要的是能够回到规格和代码，重新找到
它们，并解释某一位改变以后 CPU 或设备接下来会做什么。

\section*{怎样查 Intel 手册}

Intel Software Developer's Manual（SDM）不是一本适合从第一页顺序读到
最后一页的教材。遇到具体问题时，先在卷目录中判断问题属于哪一类，再用章节
标题、寄存器名或指令名检索。不同版次可能改变页码，章节号和标题通常比页码
更适合作为笔记索引。

\begin{osgridlongtable}{L{0.18\textwidth}L{0.31\textwidth}L{0.39\textwidth}}
\textbf{手册} & \textbf{先到这里找什么} & \textbf{本书怎样使用}\\ \hline
\endfirsthead
\textbf{手册} & \textbf{先到这里找什么} & \textbf{本书怎样使用}\\ \hline
\endhead
Volume 1 & 寄存器、数据类型、执行环境和指令集概览 &
解释汇编程序看到的寄存器、标志位和基本执行模型。\\ \hline
Volume 2 & 单条指令的编码、操作过程、标志位影响和异常 &
核对 \code{LGDT}、\code{MOV CRx}、\code{IRETQ}、
\code{SYSCALL} 等指令实际会改变什么。\\ \hline
Volume 3A--3D & 复位、保护模式、分页、中断、任务状态和系统寄存器 &
重绘复位地址、长模式切换、页表、异常现场和系统调用入口图。\\ \hline
Volume 4 & 型号相关寄存器（MSR）的编号与位字段 &
核对 EFER 等 MSR；使用前仍需通过 CPUID 判断处理器是否提供相应能力。\\ \hline
\end{osgridlongtable}

查到手册结论后，还要回到项目做三次核对。先在源码中找到写寄存器或构造数据
结构的指令；再检查 ELF、ROM 或磁盘镜像中的最终字节；最后观察 QEMU 串口
日志是否真的越过了对应状态。手册说明处理器应当怎样工作，源码说明我们希望
它怎样工作，构建产物和运行结果则用来发现两者之间的差异。
